\documentclass[10pt,a4paper]{article}

% ==================================================================
% Packages
% ==================================================================
\usepackage[T1]{fontenc}
\usepackage[utf8]{inputenc}
\usepackage{lmodern}
\usepackage[margin=2.0cm]{geometry}
\usepackage{amsmath,amssymb,amsthm}
\usepackage{mathtools}
\usepackage{bm}
\usepackage{graphicx}
\usepackage{booktabs}
\usepackage{longtable}
\usepackage{array}
\usepackage{enumitem}
\usepackage{xcolor}
\usepackage{listings}
\usepackage{titlesec}
\usepackage{fancyhdr}
\usepackage{caption}
\usepackage{hyperref}
\usepackage{tikz}
\usepackage{siunitx}
\usetikzlibrary{arrows.meta,positioning,shapes.geometric,fit,backgrounds}

\hypersetup{
  colorlinks=true,
  linkcolor=blue!50!black,
  urlcolor=blue!50!black,
  citecolor=blue!50!black,
  pdftitle={PetroLumen Reservoir Module -- Mathematical and Implementation Reference},
  pdfauthor={PetroLumen Project}
}

\pagestyle{fancy}
\fancyhf{}
\fancyhead[L]{\small PetroLumen}
\fancyhead[R]{\small Reservoir Module --- Mathematical Reference}
\fancyfoot[C]{\thepage}
\renewcommand{\headrulewidth}{0.4pt}

\definecolor{codebg}{RGB}{247,247,250}
\definecolor{codekw}{RGB}{0,90,160}
\definecolor{codestr}{RGB}{140,60,20}
\definecolor{codecom}{RGB}{90,120,90}

\lstdefinestyle{py}{
  language=Python,
  backgroundcolor=\color{codebg},
  basicstyle=\ttfamily\scriptsize,
  keywordstyle=\color{codekw}\bfseries,
  stringstyle=\color{codestr},
  commentstyle=\color{codecom}\itshape,
  numbers=left,
  numberstyle=\tiny\color{gray},
  numbersep=6pt,
  frame=single,
  rulecolor=\color{gray!40},
  breaklines=true,
  breakatwhitespace=false,
  showstringspaces=false,
  tabsize=4,
  captionpos=b
}
\lstdefinestyle{sh}{
  language=bash,
  backgroundcolor=\color{codebg},
  basicstyle=\ttfamily\scriptsize,
  keywordstyle=\color{codekw}\bfseries,
  commentstyle=\color{codecom}\itshape,
  frame=single,
  rulecolor=\color{gray!40},
  breaklines=true,
  showstringspaces=false,
  captionpos=b
}
\lstdefinestyle{txt}{
  backgroundcolor=\color{codebg},
  basicstyle=\ttfamily\scriptsize,
  frame=single,
  rulecolor=\color{gray!40},
  breaklines=true,
  showstringspaces=false,
  captionpos=b
}
\lstset{style=py}

\titleformat{\section}{\Large\bfseries\color{blue!40!black}}{\thesection}{0.7em}{}
\titleformat{\subsection}{\large\bfseries\color{blue!30!black}}{\thesubsection}{0.5em}{}
\titleformat{\subsubsection}{\normalsize\bfseries\color{blue!20!black}}{\thesubsubsection}{0.4em}{}

\newtheorem{definition}{Definition}[section]
\newtheorem{remark}{Remark}[section]

% shorthand
\newcommand{\dd}{\mathrm{d}}
\newcommand{\pd}[2]{\frac{\partial #1}{\partial #2}}
\newcommand{\grad}{\nabla}
\newcommand{\divv}{\nabla\cdot}
\newcommand{\bo}{B_o}
\newcommand{\bw}{B_w}
\newcommand{\bg}{B_g}
\newcommand{\rs}{R_s}
\newcommand{\sw}{S_w}
\newcommand{\so}{S_o}
\newcommand{\sg}{S_g}
\newcommand{\kr}{k_r}
\newcommand{\krw}{k_{rw}}
\newcommand{\kro}{k_{ro}}
\newcommand{\krg}{k_{rg}}
\newcommand{\muo}{\mu_o}
\newcommand{\muw}{\mu_w}
\newcommand{\mug}{\mu_g}
\newcommand{\pc}{p_c}
\newcommand{\phio}{\phi}
\newcommand{\ct}{c_t}
\newcommand{\inj}{\mathrm{inj}}
\newcommand{\prod}{\mathrm{prod}}
\newcommand{\res}{\mathrm{res}}

\title{
  \vspace{-1.2cm}
  \textbf{PetroLumen}\\[0.2em]
  \large Reservoir Module --- Mathematical Foundations,\\
  Numerical Methods, and Free-Tool Implementation Reference\\[0.4em]
  \normalsize From governing PDEs to validated Python code
}
\author{PetroLumen Engineering}
\date{\today}

\begin{document}
\maketitle
\thispagestyle{empty}

\begin{abstract}
\noindent
This document is the mathematical and implementation reference for the
\textbf{reservoir module} of PetroLumen. It begins with the black-oil
conservation equations and Darcy's law, derives the material-balance and
well-test equations that the module exposes as engineering tools, presents the
finite-volume discretisation and Newton--Raphson linearisation used by the
OPM Flow and JutulDarcy engines, documents the PVT correlations and
geostatistical formulations, and ends with benchmark tolerances and concrete
Python integration code. Every equation is tied to a free or open-source
implementation path so that the reservoir module can be built, validated and
shipped without commercial software licences.
\end{abstract}

\tableofcontents
\clearpage

% ==================================================================
\part{Physical and Mathematical Foundations}
% ==================================================================

\section{Governing equations of multiphase flow}

\subsection{Black-oil model: assumptions and conservation laws}

The black-oil model treats the reservoir as three phases (water, oil, gas) and
three pseudo-components (water, stock-tank oil, surface gas). Oil and gas may
exchange mass through solution gas; water is immiscible with the hydrocarbon
phases[reference:0].

Let $\alpha\in\{w,o,g\}$ denote phase, $\phi$ porosity, $\rho_\alpha$ phase
density at reservoir conditions, $S_\alpha$ phase saturation, $\bm{u}_\alpha$
Darcy velocity, and $q_\alpha$ source/sink term. The mass-conservation
equations are

\begin{align}
  \pd{}{t}\bigl(\phi\,\rho_w S_w\bigr)
  &= -\divv\bigl(\rho_w \bm{u}_w\bigr) + q_w, \label{eq:cons-w}\\[2pt]
  \pd{}{t}\bigl(\phi\,\rho_o S_o\bigr)
  &= -\divv\bigl(\rho_o \bm{u}_o\bigr) + q_o, \label{eq:cons-o}\\[2pt]
  \pd{}{t}\bigl(\phi\,\rho_g S_g\bigr)
  &= -\divv\bigl(\rho_g \bm{u}_g\bigr) + q_g. \label{eq:cons-g}
\end{align}
These are the standard black-oil conservation statements[reference:1][reference:2].

The phase velocities follow the multiphase extension of Darcy's law:
\begin{equation}
  \bm{u}_\alpha
  = -\frac{k\,\kr_\alpha(S_\alpha)}{\mu_\alpha(p_\alpha)}
    \bigl(\grad p_\alpha - \rho_\alpha\, g\,\grad z\bigr),
  \qquad \alpha=w,o,g,
  \label{eq:darcy}
\end{equation}
where $k$ is absolute permeability, $\kr_\alpha$ relative permeability,
$\mu_\alpha$ viscosity, $p_\alpha$ phase pressure, $g$ gravitational
acceleration and $z$ depth[reference:3].

The system is closed by
\begin{align}
  S_w + S_o + S_g &= 1, \label{eq:sat}\\
  p_{c,wo}(S_w) &= p_o - p_w, \label{eq:pc-wo}\\
  p_{c,go}(S_g) &= p_g - p_o, \label{eq:pc-go}
\end{align}
with capillary pressure functions $p_{c,wo}$ and $p_{c,go}$ supplied by the
rock model.

\subsection{Black-oil PVT relations}

The black-oil formulation expresses phase densities and formation volume
factors through the solution gas--oil ratio $\rs(p)$, the oil formation volume
factor $\bo(p,\rs)$, the gas formation volume factor $\bg(p)$, the water
formation volume factor $\bw(p)$, and the volatilised oil--gas ratio $R_v$.
The reservoir densities are
\begin{align}
  \rho_o &= \frac{\rho_{o,\mathrm{STC}} + \rs\,\rho_{g,\mathrm{STC}}}{\bo}, \label{eq:rho-o}\\
  \rho_g &= \frac{\rho_{g,\mathrm{STC}} + R_v\,\rho_{o,\mathrm{STC}}}{\bg}, \label{eq:rho-g}\\
  \rho_w &= \frac{\rho_{w,\mathrm{STC}}}{\bw}, \label{eq:rho-w}
\end{align}
where the subscript STC denotes standard (stock-tank) conditions.

The black-oil equations can be written in the equivalent ``$\phi S/B$'' form
used by most simulators:
\begin{equation}
  \pd{}{t}\!\left(\frac{\phi\,S_o}{\bo}\right)
  + \divv\!\left(\frac{\bm{u}_o}{\bo}\right)
  = q_o',
  \label{eq:bo-form}
\end{equation}
and analogously for water and gas, with the gas equation carrying both free-gas
and solution-gas contributions.

\subsection{Compositional extension (outline)}

For compositional runs, the black-oil pseudo-components are replaced by $N_c$
components with mole fractions $x_{i\alpha}$ in phase $\alpha$. The
conservation equations become
\begin{equation}
  \pd{}{t}\!\left(\phi\sum_{\alpha}\rho_\alpha\,S_\alpha\,x_{i\alpha}\right)
  + \divv\!\left(\sum_{\alpha}\rho_\alpha\,x_{i\alpha}\,\bm{u}_\alpha\right)
  = q_i,
  \qquad i=1,\dots,N_c,
  \label{eq:comp-cons}
\end{equation}
with phase equilibrium computed from an equation of state (e.g.\ Peng--Robinson)
and fugacity equality. OPM Flow and JutulDarcy both support compositional
formulations; PetroLumen exposes them through the same engine plug-in
interface but with a different fluid model type.

% ------------------------------------------------------------------
\section{Volumetrics and original fluids in place}
% ------------------------------------------------------------------

\subsection{Volumetric OOIP and OGIP}

For a grid cell $i$ with bulk volume $V_{b,i}$, porosity $\phi_i$, and initial
water saturation $S_{w,i}^0$, the original oil in place is
\begin{equation}
  N_i = \frac{V_{b,i}\,\phi_i\,(1-S_{w,i}^0)}{\bo_i^0},
  \label{eq:ooip-cell}
\end{equation}
and the original gas in place (free gas) is
\begin{equation}
  G_i = \frac{V_{b,i}\,\phi_i\,S_{g,i}^0}{\bg_i^0}.
  \label{eq:ogip-cell}
\end{equation}
Summing over active cells gives the field totals.

\subsection{Recovery factor and drive indices}

The recovery factor at time $t$ is
\begin{equation}
  \mathrm{RF}(t) = \frac{N_p(t)}{N},
  \label{eq:rf}
\end{equation}
where $N_p$ is cumulative oil production.

The drive indices of Pirson partition the total withdrawal into expansion
terms:
\begin{align}
  \mathrm{DI}_{\mathrm{oil}} &= \frac{N(\bo-\bo_i)}{\bo_i}, \\
  \mathrm{DI}_{\mathrm{gas}} &= \frac{N\,m\,\bo_i\!\left(\frac{\bg}{\bg_i}-1\right)}{\bo_i}, \\
  \mathrm{DI}_{\mathrm{water}} &= \frac{W_e - W_p\,\bw}{\bo_i}, \\
  \mathrm{DI}_{\mathrm{rock}} &= \frac{N\,\bo_i\!\left(\frac{c_w S_w + c_f}{1-S_w}\right)\Delta p}{\bo_i},
\end{align}
where $m$ is the initial gas-cap-to-oil-zone volume ratio, $W_e$ water influx,
$W_p$ water production, $c_w$ water compressibility and $c_f$ formation
compressibility. The sum of drive indices equals one at any time.

% ------------------------------------------------------------------
\section{Material balance}
% ------------------------------------------------------------------

\subsection{General Havlena--Odeh form}

The general material-balance equation for an oil reservoir is written as
\begin{equation}
  F = N\bigl(E_o + m\,E_g + E_{fw}\bigr) + W_e\,\bw,
  \label{eq:mbe}
\end{equation}
where
\begin{align}
  F &= N_p\!\left[\bo + (R_p-\rs)\bg\right] + W_p\,\bw, \label{eq:F}\\
  E_o &= (\bo-\bo_i) + (R_{si}-\rs)\bg, \label{eq:Eo}\\
  E_g &= \bo_i\!\left(\frac{\bg}{\bg_i}-1\right), \label{eq:Eg}\\
  E_{fw} &= \bo_i\!\left[\frac{c_w S_{wi}+c_f}{1-S_{wi}}\right]\Delta p. \label{eq:Efw}
\end{align}
$F$ is the underground withdrawal, $E_o$ the oil-and-dissolved-gas expansion,
$E_g$ the gas-cap expansion, and $E_{fw}$ the connate-water and rock
expansion[reference:4][reference:5].

\subsection{Linearised Havlena--Odeh plots}

For a reservoir with no initial gas cap ($m=0$) and no water influx, equation
\eqref{eq:mbe} reduces to
\begin{equation}
  F = N\,(E_o + E_{fw}),
  \label{eq:mbe-simple}
\end{equation}
so a plot of $F$ against $E_o+E_{fw}$ yields a straight line through the origin
with slope $N$.

When a gas cap is present, the linearised form is
\begin{equation}
  \frac{F - W_e\bw}{E_o + \bo_i E_{fw}}
  = N + G\left(\frac{E_g + \bg_i E_{fw}}{E_o + \bo_i E_{fw}}\right),
  \label{eq:havlena-odeh}
\end{equation}
where $G$ is the original free gas in place. A plot of the left-hand side
against the bracket on the right gives intercept $N$ and slope $G$[reference:6].

\subsection{Gas material balance}

For a volumetric dry-gas reservoir, the material balance reduces to the
$p/Z$ straight line:
\begin{equation}
  \frac{p}{Z} = \frac{p_i}{Z_i}\left(1 - \frac{G_p}{G}\right).
  \label{eq:pz}
\end{equation}
A linear regression of $p/Z$ against $G_p$ gives $G$ as the $x$-intercept
[reference:7].

\subsection{Aquifer influx}

The Van Everdingen--Hurst dimensionless aquifer influx is
\begin{equation}
  W_e(t) = B\,\Delta p\,W_D(t_D),
  \label{eq:veh}
\end{equation}
with dimensionless time
\begin{equation}
  t_D = \frac{k\,t}{\phi\,\mu\,c_t\,r_a^2}.
  \label{eq:td-aquifer}
\end{equation}
Carter--Tracy provides a computationally efficient alternative; PetroLumen
exposes both through a single \texttt{AquiferModel} interface, with
ECLIPSE \texttt{AQUTAB} export[reference:8].

% ------------------------------------------------------------------
\section{Pressure transient analysis}
% ------------------------------------------------------------------

\subsection{Dimensionless variables}

Dimensionless pressure and time for a single-phase oil well are
\begin{align}
  p_D &= \frac{2\pi k h}{q B \mu}\,\Delta p, \label{eq:pd}\\
  t_D &= \frac{k\,\Delta t}{\phi\,\mu\,c_t\,r_w^2}. \label{eq:td}
\end{align}
These definitions underpin type-curve matching and semi-log analysis
[reference:9].

\subsection{Radial-flow semi-log analysis}

During infinite-acting radial flow, the drawdown or buildup pressure follows
\begin{equation}
  \Delta p = \frac{162.6\,q B \mu}{k h}\,
  \left[\log_{10}\!\left(\frac{k t}{\phi \mu c_t r_w^2}\right) - 3.23 + 0.87\,s\right],
  \label{eq:semilog}
\end{equation}
where $s$ is the skin factor. The slope $m$ of the semi-log straight line
(psi per log$_{10}$ cycle) gives the permeability:
\begin{equation}
  k = \frac{162.6\,q B \mu}{m h}.
  \label{eq:k-semilog}
\end{equation}
The skin is obtained from the pressure change at one hour:
\begin{equation}
  s = 1.151\left[\frac{\Delta p_{1\mathrm{hr}}}{m}
  - \log_{10}\!\left(\frac{k}{\phi \mu c_t r_w^2}\right) + 3.23\right].
  \label{eq:skin}
\end{equation}
These are the standard equations used in commercial PTA software and reproduced
by \texttt{welltest-pta}[reference:10].

\subsection{Horner buildup analysis}

For a buildup test, the Horner equation is
\begin{equation}
  p_{ws} = p_i - \frac{162.6\,q B \mu}{k h}
  \log_{10}\!\left(\frac{t_p + \Delta t}{\Delta t}\right),
  \label{eq:horner}
\end{equation}
where $t_p$ is the producing time before shut-in and $\Delta t$ is shut-in
time[reference:11][reference:12]. A plot of $p_{ws}$ against the Horner time
ratio yields a straight line whose slope gives $k$ and whose extrapolation to
unity gives $p_i$ (or the current average pressure).

\subsection{Bourdet derivative}

The Bourdet logarithmic derivative is
\begin{equation}
  p' = \frac{\dd p}{\dd \ln t},
  \label{eq:bourdet}
\end{equation}
computed with a three-point window to smooth noise while preserving flow-regime
signatures. \texttt{welltest-pta} exposes \texttt{bourdet\_derivative(dt, dp, L=0.2, use\_smooth=True)}
and performs automatic flow-regime identification by fitting piece-wise
log--log slopes to the derivative curve[reference:13].

% ------------------------------------------------------------------
\section{Geostatistics and static modelling}
% ------------------------------------------------------------------

\subsection{Variogram}

The experimental semi-variogram is
\begin{equation}
  \hat{\gamma}(\bm{h}) = \frac{1}{2N(\bm{h})}
  \sum_{i=1}^{N(\bm{h})}\bigl[z(\bm{x}_i)-z(\bm{x}_i+\bm{h})\bigr]^2,
  \label{eq:variogram}
\end{equation}
where $\bm{h}$ is the lag vector and $N(\bm{h})$ the number of pairs at that
lag. Permissible models (spherical, exponential, Gaussian) are fitted to
$\hat{\gamma}$; \texttt{GeostatsPy} provides \texttt{gam\_3D},
\texttt{make\_variogram\_3D} and \texttt{vmodel\_3D}[reference:14].

\subsection{Ordinary kriging}

The ordinary kriging estimator is
\begin{equation}
  \hat{z}(\bm{x}_0) = \sum_{i=1}^{n}\lambda_i\,z(\bm{x}_i),
  \label{eq:ok-est}
\end{equation}
with weights $\lambda_i$ solving the kriging system
\begin{equation}
  \begin{bmatrix}
    \bm{\Gamma} & \bm{1}\\
    \bm{1}^{\!\top} & 0
  \end{bmatrix}
  \begin{bmatrix}
    \bm{\lambda}\\ \mu
  \end{bmatrix}
  =
  \begin{bmatrix}
    \bm{\gamma}_0\\ 1
  \end{bmatrix},
  \label{eq:ok-system}
\end{equation}
where $\Gamma_{ij}=\gamma(\bm{x}_i-\bm{x}_j)$, $(\bm{\gamma}_0)_i=\gamma(\bm{x}_i-\bm{x}_0)$,
and $\mu$ is the Lagrange multiplier. \texttt{GeostatsPy} implements this as
\texttt{kb2d}/\texttt{kb3d}[reference:15].

\subsection{Sequential Gaussian simulation}

SGS proceeds by (i) normal-score transforming the data, (ii) selecting a random
path through the grid, (iii) kriging the mean and variance at each unsampled
node from previously simulated and original data, and (iv) drawing a value from
the local Gaussian distribution and adding it to the simulated field:
\begin{equation}
  z^{*}(\bm{x}_0) = \hat{z}(\bm{x}_0) + \sigma_K(\bm{x}_0)\,r,
  \qquad r\sim\mathcal{N}(0,1),
  \label{eq:sgs}
\end{equation}
where $\sigma_K$ is the kriging standard deviation. \texttt{GeostatsPy}
provides \texttt{sgsim\_3D} and the associated normal-score utilities
[reference:16].

% ==================================================================
\part{Numerical Methods}
% ==================================================================

\section{Finite-volume discretisation}

\subsection{Two-point flux approximation}

OPM Flow and JutulDarcy use a finite-volume discretisation with the two-point
flux approximation (TPFA) on corner-point or unstructured grids[reference:17]
[reference:18]. For a phase $\alpha$, the flux across the face between cells
$i$ and $j$ is
\begin{equation}
  F_{\alpha,ij} = -T_{ij}\,\lambda_{\alpha,ij}\,
  \bigl(p_{\alpha,j}-p_{\alpha,i} - \rho_{\alpha,ij}\,g\,(z_j-z_i)\bigr),
  \label{eq:tpfa}
\end{equation}
where $T_{ij}$ is the transmissibility,
\begin{equation}
  T_{ij} = \frac{A_{ij}\,k_{ij}}{d_{ij}},
  \label{eq:trans}
\end{equation}
with $A_{ij}$ the interfacial area, $k_{ij}$ the harmonic-mean permeability,
and $d_{ij}$ the centroid-to-centroid distance. The phase mobility is
\begin{equation}
  \lambda_{\alpha,ij} = \left(\frac{\kr_\alpha}{\mu_\alpha B_\alpha}\right)_{\!ij},
  \label{eq:mobility}
\end{equation}
evaluated upstream.

\subsection{Time discretisation}

The backward-Euler (fully implicit) discretisation of the mass-conservation
equation for cell $i$ and phase $\alpha$ is
\begin{equation}
  \frac{V_i}{\Delta t}
  \Bigl[(\phi\,\rho_\alpha S_\alpha)_i^{n+1}
       -(\phi\,\rho_\alpha S_\alpha)_i^{n}\Bigr]
  + \sum_{j\in\mathcal{N}(i)} F_{\alpha,ij}^{n+1}
  = q_{\alpha,i}^{n+1},
  \label{eq:be}
\end{equation}
where $\mathcal{N}(i)$ is the set of neighbouring cells[reference:19].

\subsection{IMPES versus fully implicit}

IMPES (implicit pressure, explicit saturation) solves a single pressure equation
implicitly and updates saturations explicitly; it is cheap per timestep but
conditionally stable and unsuitable for large timesteps or strong
capillary/gravity effects. Fully implicit (FIM) solves all conservation
equations simultaneously at $t^{n+1}$, requiring a Newton--Raphson iteration
but remaining stable for large timesteps. OPM Flow's default is FIM; IMPES is
available for screening and teaching[reference:20].

% ------------------------------------------------------------------
\section{Newton--Raphson linearisation}
% ------------------------------------------------------------------

Let $\bm{R}(\bm{x})=0$ be the discrete residual vector, with $\bm{x}$ the vector
of primary unknowns (typically $p_o$, $S_w$, $S_g$). Newton--Raphson iterates
\begin{equation}
  \bm{J}(\bm{x}^k)\,\delta\bm{x}^k = -\bm{R}(\bm{x}^k),
  \qquad
  \bm{x}^{k+1} = \bm{x}^k + \delta\bm{x}^k,
  \label{eq:newton}
\end{equation}
where $\bm{J}=\partial\bm{R}/\partial\bm{x}$ is the Jacobian. Each Newton step
requires the solution of a large, sparse, ill-conditioned linear system
[reference:21][reference:22].

The Jacobian has a block structure: diagonal blocks couple the unknowns within a
cell, and off-diagonal blocks couple neighbouring cells through the TPFA flux
terms. For a three-phase black-oil model with three unknowns per active cell,
the block size is $3\times 3$.

\subsection{Convergence criteria}

OPM Flow and JutulDarcy typically require
\begin{align}
  \|\bm{R}\|_2 &< \varepsilon_R, \label{eq:conv-res}\\
  \|\delta\bm{x}\|_2 &< \varepsilon_x, \label{eq:conv-dx}
\end{align}
with tolerances scaled by the number of active cells and the magnitude of the
solution. Typical defaults are $\varepsilon_R\sim10^{-6}$ and
$\varepsilon_x\sim10^{-6}$ in consistent units; PetroLumen should expose these
as solver options and log the achieved residuals per Newton iteration.

\subsection{Linear solvers}

The linear system $\bm{J}\delta\bm{x}=-\bm{R}$ is solved with a Krylov method
(GMRES or BiCGStab) preconditioned by CPR (constrained pressure residual) or
AMG. CPR splits the system into a pressure part solved with AMG and a
transport part solved with ILU, which is the standard strategy in OPM Flow
[reference:23][reference:24]. GPU-accelerated CPR variants are documented for
SPE10 and Norne[reference:25].

% ------------------------------------------------------------------
\section{Well modelling}
% ------------------------------------------------------------------

Wells are coupled to the reservoir equations through a well model. For a well
$w$ completed in cells $\mathcal{C}_w$, the well equations are solved fully
implicitly with the reservoir equations[reference:26]:
\begin{align}
  q_{\alpha,w} &= \sum_{c\in\mathcal{C}_w} W_{w,c}\,\lambda_{\alpha,c}\,
  \bigl(p_{\alpha,c} - p_{w} - \rho_{\alpha,c}\,g\,(z_c-z_w)\bigr), \label{eq:well-rate}\\
  \sum_{\alpha} q_{\alpha,w} &= q_{w}^{\mathrm{target}}, \label{eq:well-constraint}
\end{align}
where $W_{w,c}$ is the well index, $p_w$ the wellbore pressure, and
$q_w^{\mathrm{target}}$ the rate or pressure constraint. The well index for a
vertical well in a rectangular cell is
\begin{equation}
  W_{w,c} = \frac{2\pi\,k_c\,h_c}
  {\ln(r_{e,c}/r_w) + s_c},
  \label{eq:well-index}
\end{equation}
with $r_{e,c}$ the equivalent radius and $s_c$ the skin.

% ==================================================================
\part{Free-Tool Implementation}
% ==================================================================

\section{Engine plug-in interface}

All simulation engines implement the same contract:

\begin{lstlisting}[caption={Engine plug-in protocol.}]
from typing import Protocol, runtime_checkable

@runtime_checkable
class EnginePlugin(Protocol):
    name: str
    def validate_case(self, case) -> list[str]: ...
    def build_inputs(self, case, workdir) -> dict: ...
    def run(self, workdir, options) -> "RunResult": ...
    def parse_results(self, workdir) -> "ResultSet": ...
\end{lstlisting}

\subsection{OPM Flow back-end}

OPM Flow reads Eclipse-style decks. A minimal writer:

\begin{lstlisting}[caption={Deck writer for a black-oil case.}]
def write_deck(case, grid_path, out_path):
    with open(out_path, "w") as f:
        f.write("RUNSPEC\n")
        f.write("DIMENS\n 10 10 3 /\n")
        f.write("OIL\nWATER\nGAS\nFIELD\n")
        f.write("GRID\n")
        f.write(f"INCLUDE\n '{grid_path}' /\n")
        f.write("PROPS\n")
        f.write(f"INCLUDE\n '{case.fluid_id}.inc' /\n")
        f.write(f"INCLUDE\n '{case.rock_id}.inc' /\n")
        f.write("SCHEDULE\n")
        f.write(f"INCLUDE\n '{case.schedule_id}.inc' /\n")
        f.write("END\n")
\end{lstlisting}

Run with subprocess isolation to keep the GPL boundary clean:

\begin{lstlisting}[caption={Subprocess wrapper for OPM Flow.}]
import subprocess, hashlib, datetime

def run_flow(deck_path, workdir, timeout=86400):
    proc = subprocess.run(
        ["flow", deck_path], cwd=workdir,
        capture_output=True, text=True, timeout=timeout
    )
    return {
        "deck_hash": hashlib.sha256(open(deck_path, "rb").read()).hexdigest(),
        "exit_code": proc.returncode,
        "finished": datetime.datetime.utcnow().isoformat(),
    }
\end{lstlisting}

\subsection{JutulDarcy back-end}

JutulDarcy is a Julia package and is fully differentiable with respect to
forces and discretisation parameters. Call it from Python via
\texttt{juliacall} or a Julia subprocess. Its differentiability makes it
particularly useful for gradient-based history matching and optimisation.

\subsection{pyMRST back-end}

pyMRST wraps the MATLAB Reservoir Simulation Toolbox and is GPL-3.0. Use it for
research-grade discretisations and teaching; like OPM Flow, invoke it as an
external process.

% ------------------------------------------------------------------
\section{PVT implementation}
% ------------------------------------------------------------------

\subsection{Standing correlation}

Standing's bubble-point pressure is
\begin{equation}
  p_b = 18.2\left[\left(\frac{\rs}{\gamma_g}\right)^{0.83}10^{a}-1\right],
  \qquad
  a = 0.00091\,T - 0.0125\,\mathrm{API},
  \label{eq:standing}
\end{equation}
with $p_b$ in psia, $\rs$ in scf/STB, $T$ in \si{\fahrenheit} and API gravity
in degrees[reference:28].

\subsection{Vasquez--Beggs correlation}

Vasquez and Beggs give
\begin{equation}
  \rs = C_1\,\gamma_{gs}\,p^{C_2}\exp\!\left(\frac{C_3\,\mathrm{API}}{T+460}\right),
  \label{eq:vb-rs}
\end{equation}
with coefficients selected for API $\le 30$ or $>30$[reference:29][reference:30].
The oil formation volume factor is
\begin{equation}
  \bo = 1.0 + C_1\rs + C_2(T-60)\left(\frac{\mathrm{API}}{\gamma_{gs}}\right)
  + C_3\rs(T-60)\left(\frac{\mathrm{API}}{\gamma_{gs}}\right).
  \label{eq:vb-bo}
\end{equation}

\subsection{Hall--Yarborough Z-factor}

The Hall--Yarborough equation of state is used for gas $Z$-factor:
\begin{equation}
  Z = \frac{0.06125\,p_{pr}}{Y}\exp\!\bigl[-1.2(1-1/T_{pr})^2\bigr],
  \label{eq:hy-z}
\end{equation}
where $Y$ solves the implicit equation
\begin{equation}
  -0.06125\,p_{pr}\frac{Y}{T_{pr}}
  + \frac{Y+Y^2+Y^3-Y^4}{(1-Y)^3}
  - \left(\frac{14.76}{T_{pr}}-\frac{9.76}{T_{pr}^2}+\frac{4.58}{T_{pr}^3}\right)Y^2
  + \left(\frac{90.7}{T_{pr}}-\frac{242.2}{T_{pr}^2}+\frac{42.4}{T_{pr}^3}\right)Y^{(2.18+2.82/T_{pr})}
  = 0.
  \label{eq:hy-implicit}
\end{equation}
Solve by Newton iteration; \texttt{pyResToolbox} provides a tested
implementation.

% ------------------------------------------------------------------
\section{Well-test implementation}
% ------------------------------------------------------------------

\texttt{welltest-pta} provides the full pipeline:
\begin{lstlisting}[caption={Well-test analysis with welltest-pta.}]
from welltest_pta import Event
event = Event(dt, dp, rate_history)
dt_b, dp_b = event.bourdet(L=0.2, use_smooth=True)
params = event.reservoir_params(k=..., h=..., mu=..., B=...)
\end{lstlisting}
The library exposes \texttt{.bourdet()}, \texttt{.horner()}, \texttt{.mdh()} and
\texttt{.flow\_regimes()} on every drawdown/buildup event[reference:31].

% ------------------------------------------------------------------
\section{Material-balance implementation}
% ------------------------------------------------------------------

\texttt{PetroCalc} provides \texttt{petrocalc.reservoir} and
\texttt{petrocalc.gas\_reservoir} with material-balance, drive-index and
recovery functions[reference:32]. The Havlena--Odeh regression can be assembled
directly with \texttt{NumPy}:

\begin{lstlisting}[caption={Havlena--Odeh linear regression.}]
import numpy as np

def havlena_odeh(F, Eo, Eg, Ef, We, Bw):
    y = (F - We * Bw) / (Eo + Ef)
    x = (Eg + Ef) / (Eo + Ef)
    A = np.vstack([np.ones_like(x), x]).T
    (N, G), *_ = np.linalg.lstsq(A, y, rcond=None)
    return N, G
\end{lstlisting}

% ------------------------------------------------------------------
\section{Geostatistics implementation}
% ------------------------------------------------------------------

\begin{lstlisting}[caption={Variogram modelling and kriging with GeostatsPy.}]
import geostatspy.geostats as geostats
import geostatspy.GSLIB as GSLIB

vario = GSLIB.make_variogram(nug=0.0, nst=1, it1=1,
                            cc1=1.0, a1=1000.0)
est, var = geostats.kb2d(df, xcol="X", ycol="Y", vcol="Por",
                         nx=50, xmn=0.0, xsiz=25.0,
                         nz=1, zmin=0.0, zsiz=1.0,
                         ndmin=1, ndmax=8, nmult=1,
                         ktype=1, vario=vario)
\end{lstlisting}
\texttt{geostats.sgsim\_3D} performs sequential Gaussian simulation
[reference:33][reference:34].

% ------------------------------------------------------------------
\section{RESQML and grid I/O}
% ------------------------------------------------------------------

\texttt{resqpy} is a pure Python package for reading, writing and modifying
RESQML models; it implements the full RESQML units-of-measure system
[reference:35][reference:36].

\begin{lstlisting}[caption={Opening a RESQML model and accessing a grid.}]
import resqpy.model as rq

model = rq.Model("reservoir.epc")
grid = model.grid()          # first grid in the model
extent = grid.extent_kji     # (nk, nj, ni)
\end{lstlisting}

\texttt{Petres} constructs corner-point grids programmatically and round-trips
\texttt{.GRDECL} files.

% ==================================================================
\part{Validation}
% ==================================================================

\section{Benchmark tolerances}

\begin{longtable}{@{}p{2.0cm}p{5.4cm}p{6.8cm}@{}}
\toprule
\textbf{Benchmark} & \textbf{Physics} & \textbf{Acceptance criterion}\\
\midrule
\endhead
SPE1  & Three-phase, two wells, gas coning & BHP and gas rate within 2\% of published reference\\
SPE9  & Black-oil, 26 wells, faulted grid & Cumulative production within 2\% at final time\\
SPE10 & 60$\times$220$\times$85 heterogeneous permeability & Water-cut history within 3\%; determinism across runs\\
Norne & Real field, complex geology & Summary vectors within agreed tolerance vs.\ reference\\
\bottomrule
\end{longtable}

\noindent Additional gates:
\begin{itemize}[leftmargin=1.4em]
  \item \textbf{Determinism.} Same deck and engine version must produce
        bit-identical summary output on the same hardware.
  \item \textbf{Unit round-trip.} SI $\rightarrow$ field $\rightarrow$ SI
        must be lossless within \num{1e-12} relative error.
  \item \textbf{Deck round-trip.} Parse an Eclipse deck, regenerate it, and
        confirm OPM Flow produces the same summary vectors.
  \item \textbf{RESQML round-trip.} RESQML in $\rightarrow$ internal $\rightarrow$
        RESQML out must be lossless for supported attributes.
\end{itemize}

% ------------------------------------------------------------------
\section{Regression harness}

\begin{lstlisting}[caption={Pytest fixture for SPE1 regression.}]
import pytest, subprocess, json

@pytest.mark.benchmark
def test_spe1(tmp_path):
    deck = tmp_path / "SPE1.DATA"
    deck.write_text(SPE1_DECK)
    result = subprocess.run(["flow", str(deck)], cwd=tmp_path,
                            capture_output=True, text=True, timeout=600)
    assert result.returncode == 0
    summary = parse_summary(tmp_path / "SPE1.SMSPEC")
    assert abs(summary["FOPT"][-1] - 1.0e6) / 1.0e6 < 0.02
\end{lstlisting}

% ==================================================================
\part{Deployment and Engineering Practice}
% ==================================================================

\section{Licensing boundaries}

\begin{longtable}{@{}p{3.4cm}p{2.4cm}p{8.4cm}@{}}
\toprule
\textbf{Component} & \textbf{Licence} & \textbf{Integration rule}\\
\midrule
OPM Flow & OSL-3.0 / GPL & External executable; never link into PetroLumen\\
JutulDarcy.jl & MIT & May be imported via \texttt{juliacall}, but ship as optional extra\\
pyMRST & GPL-3.0 & External process only\\
pyResToolbox & GPL-3.0 & External service or optional extra\\
resqpy & MIT & Safe to import directly\\
Petres & LGPL & Safe to import; dynamic linking required for LGPL compliance\\
GeostatsPy & MIT & Safe to import directly\\
welltest-pta & MIT & Safe to import directly\\
PetroCalc & MIT & Safe to import directly\\
\bottomrule
\end{longtable}

Ship PetroLumen itself under Apache-2.0 or MIT and declare GPL engines as
optional, user-installed back-ends discovered at runtime.

\section{Performance strategy}

\begin{enumerate}[leftmargin=1.4em]
  \item \textbf{Delegate numerics.} OPM Flow, JutulDarcy and PETSc carry the
        linear algebra; Python orchestrates.
  \item \textbf{Vectorise the glue.} NumPy and \texttt{xarray} for result
        aggregation; never loop over grid cells in Python.
  \item \textbf{Compile hotspots.} Numba or Cython for custom upscaling and
        derivative computation.
  \item \textbf{Parallelism by job.} Run cases concurrently on the queue; let
        the engine manage MPI/OpenMP inside a case.
  \item \textbf{Optional GPU path.} JutulDarcy on CUDA-capable hardware; keep
        it optional so the default install stays portable.
\end{enumerate}

\section{Roadmap}

\begin{longtable}{@{}p{1.4cm}p{4.6cm}p{8.2cm}@{}}
\toprule
\textbf{Phase} & \textbf{Deliverable} & \textbf{Mathematical scope}\\
\midrule
M0 & Skeleton & Schema, units, project store, CI\\
M1 & Engineering basics & Volumetrics \eqref{eq:ooip-cell}--\eqref{eq:ogip-cell},
material balance \eqref{eq:mbe}, semi-log PTA \eqref{eq:k-semilog}--\eqref{eq:skin}\\
M2 & Simulation MVP & TPFA \eqref{eq:tpfa}, backward Euler \eqref{eq:be},
Newton \eqref{eq:newton}, OPM Flow back-end, SPE1/SPE9 passing\\
M3 & Static link & Variogram \eqref{eq:variogram}, kriging \eqref{eq:ok-system},
SGS \eqref{eq:sgs}, RESQML via resqpy, SPE10 passing\\
M4 & Post-processing & Summary parsing, 3D visualisation, forecasting\\
M5 & Uncertainty & Ensembles, sensitivity, surrogate-assisted HM, Norne passing\\
M6 & Hardening & Performance tuning, plugin SDK, container image\\
\bottomrule
\end{longtable}

\section*{References}
\addcontentsline{toc}{section}{References}
\begin{enumerate}[leftmargin=1.6em,label={[\arabic*]}]
  \item L.P. Dake, \emph{Fundamentals of Reservoir Engineering}, Elsevier, 1978.
  \item T. Ahmed, \emph{Reservoir Engineering Handbook}, Gulf Professional Publishing, 5th ed.\ 2019.
  \item K.-A. Lie, \emph{An Introduction to Reservoir Simulation Using MATLAB/GNU Octave}, Cambridge University Press, 2019.
  \item D.W. Peaceman, \emph{Fundamentals of Numerical Reservoir Simulation}, Elsevier, 1977.
  \item M.J. Pyrcz, C.V. Deutsch, \emph{Geostatistical Reservoir Modeling}, Oxford University Press, 2nd ed.\ 2014.
  \item R.N. Horne, \emph{Modern Well Test Analysis: A Computer-Aided Approach}, Petroway, 1995.
  \item Open Porous Media Initiative, \emph{OPM Flow} documentation.
  \item SINTEF, \emph{JutulDarcy.jl} documentation.
  \item bp, \emph{resqpy} documentation (RESQML in Python).
  \item \emph{GeostatsPy} documentation and GSLIB translations.
  \item \emph{welltest-pta} documentation.
  \item \emph{PetroCalc} documentation.
  \item SPE Comparative Solution Projects SPE1, SPE9, SPE10; Norne benchmark dataset.
\end{enumerate}

\end{document}